\chapter{Conclusion}

    \paragraph{}
    This first stage of the thesis has allowed us to wander with deliberate care through the intellectual landscape that binds together ontology engineering, knowledge graphs, natural language processing, reinforcement learning and multi agent orchestration, revealing not only the conceptual foundations of our approach but also the delicate interdependencies that govern how regulatory knowledge may be faithfully extracted and thoughtfully extended within a modern systems engineering environment. By examining the literature with patience and precision, we now understand that knowledge extraction and knowledge extension are not opposing poles but complementary movements in a single methodological arc, each offering strengths that the other cannot supply. Extraction offers the sobriety of structure, grounding regulatory meaning in formal ontologies, curated triples and verifiable provenance, while extension brings the capacity for interpretative depth, adaptive reasoning and contextual intelligence, allowing the system to propose new relationships, resolve ambiguities and prioritise obligations as design contexts evolve. Through this review we have come to appreciate that the true promise of our framework lies not solely in the rigour of extraction nor solely in the creativity of extension but in the quiet harmony that emerges when the two are allowed to inform and refine each other, shaping a system capable of guiding engineers through the labyrinth of requirements with both clarity and grace.

    \paragraph{}
    Although this thesis has thus far remained within the contemplative domain of theoretical study, the pathway ahead is now much more clearly illuminated, and it is with this strengthened understanding that we look towards Thesis B and Thesis C, where these ideas will be translated into a living computational architecture and then evaluated with the depth and care that Sunswift’s engineering eco-system demands. 

    \newpage
    \paragraph{} 
    The implementation phase will challenge us to instantiate the ontology, construct the knowledge graph, build the extraction pipelines, develop the inference and recommender agents and integrate them through the Model Context Protocol, while the evaluation phase will require us to test the system with both synthetic and real regulatory scenarios, ensuring that its outputs remain trustworthy, explainable and aligned with the Australian Design Rules. Having completed a literature review that spans foundational theory, applied research and hybrid methods across industries as varied as aerospace, pharmaceuticals, additive manufacturing and digital twins, we now stand at the threshold of the design process with a richer intellectual compass, equipped not only with methods but with an understanding of how those methods must be shaped, constrained and harmonised if they are to serve the needs of a complex automotive design environment. In this sense, Thesis A concludes not with finality but with a sense of poised anticipation, for the ideas assembled here form the quiet architecture from which the system will grow, evolve and ultimately demonstrate its value in the chapters yet to come.